\chapter*{IV.17.15　光滑型\protect\\增广与初等奇点}
\addcontentsline{toc}{chapter}{IV.17.15　光滑型：增广与初等奇点}
\markboth{IV.17.15　光滑型：增广与初等奇点}{IV.17.15　光滑型：增广与初等奇点}
\begingroup
\def\OO{\mathcal O}\def\m{\mathfrak m}\def\mm{\mathfrak m}\def\Sym{\operatorname{Sym}}
\noindent\textbf{题名说明。}\footnote{“第17节”是英译编辑层的节题。法文第112页在“光滑型”后原题有划改，手写补题为“增广与初等奇点”，与\href{https://webusers.imj-prg.fr/~leila.schneps/grothendieckcircle/EGA/EGAVcontents.pdf}{法文扫描目录}相符。本页是二次型引言，后续15.1才是一般型，不把本节题目限为二次。英译将一部分标为难辨，并另有辨读不定的“新 EGA”字样，法文该标题处不据此补入。}

\medskip
\noindent\textbf{注意。} 我刚发现，我在关于超奇异性的 SUP 稿中引入的术语不合理，尤其与通行术语冲突。不管怎样，你大概已经注意到，那篇稿子定义1(c)中应当读作“退化”，而不是“非退化”。

问题出在特征$2$下的偶数维簇，例如以奇数维环境簇中的一个方程$\varphi=0$所定义者。按我笔记的术语，这样的簇不可能有“普通奇点”（即普通二次奇点），也就是说，$\varphi$不可能有“普通奇异零点”，因为特征$2$下奇数个变量的二次型总是“退化”的。然而，即使在特征$2$，人们一向也把仿射锥$x^2+yz=0$的原点看作普通二次奇点。\footnote{法文此处用带戏谑意味的“canular”；英译者加注，指出这也是巴黎高等师范学校的行话。正文按上下文译为“问题”。}

在本笔记中，我提出有限型局部自由模上的\emph{光滑}（或“普通”）二次型这一概念，使得
\[
\text{非退化}\ \Longrightarrow\ \text{光滑}.
\]
当$E$的秩为偶数，或者基底$S$的所有剩余特征都不等于$2$时，逆命题也成立。

有了这个概念，我提议保留 SUP 稿中的“超奇异”一词（它不与任何通行术语冲突），用来对应不能由非退化二次型描述的情形。\footnote{法112字面写“forme quadratique non dégénérée”，英74也印作“a non-degenerate quadratic form”；两者连接到“超奇异”时都缺少否定。紧下句明确写“核不能由一个非退化二次型生成”，而后句又将能由非退化二次型描述的奇异零点命名为“非超奇异”。原稿开头对定义1(c)的订正也指向同一极性，故正文依上下文纠正。表述取“不能由非退化二次型描述”，不是要求核必由某个退化二次型生成。}
相应地，在同一意义下，也可谈局部诺特预概形的超奇异点，或域上局部有限型预概形的几何超奇异点：考察局部环$A$的同态
\[
\operatorname{Sym}(\mm/\mm^2)\longrightarrow
\operatorname{gr}_{\mm}(A),
\]
要求其核不能由一个非退化二次型生成。因此，SUP 中原先误称为“普通奇异零点”的，应称为“非超奇异的奇异零点”；若想沿用接近特征$0$通行用语的名称，也可称为“非退化二次奇异零点”。类似地，对局部诺特预概形可说“非退化二次奇点”；若预概形在域$k$上局部有限型，则还可说“几何非退化二次奇点”。\footnote{法113开头明确给出局部诺特与域上局部有限型的分别；英译压缩了这一区别。SUP 是前一稿件的标签，不是集合并号。}

与此相对，“普通奇异零点”（更好是“普通二次零点”），以及“普通二次奇点”“几何普通二次奇点”，应按通常用法理解为对应光滑二次型的概念。此外，或可用“初等”代替“普通”，并把这个术语推广到不一定是二次、而可以具有任意重数的奇点。请告诉我你对此的看法。\footnote{英译者在此将受话人标为迪厄多内；法文仅写第二人称。}

另一方面，SUP 的文字看来在形式上是正确的；特别是命题7，实际引入的是超奇异零点。尽管如此，至少也还应在备注中引入$V(\varphi)^{\mathrm{sup\,sing}}$的子预概形
\[
V(\varphi)^{\mathrm{ult\,sing}}
\subseteq V(\varphi)^{\mathrm{sup\,sing}}.
\]
它对应极超奇异零点，即奇异却不是初等二次的那些零点。其描述方法大体与$V(\varphi)^{\mathrm{sup\,sing}}$相同，只是用下面将引入的“修正判别式”的零点预概形，代替通常判别式的零点预概形。因此，SUP 看来只须作术语上的修订（当然，也应改动 SUP 第10页引入的相关术语）。
\paragraph{译校说明（非原文）。}
\paragraph{光滑的对象与极化根基。}
按本段的用法，“光滑二次型”说的是射影零点二次超曲面的光滑性，不是说仿射锥也在顶点光滑。设$q$为域上的二次型，其极化为
\[
b_q(v,w)=q(v+w)-q(v)-q(w).
\]
在任意特征，$q$在$v$处的一阶微分就是$b_q(v,-)$。因此射影二次超曲面的几何奇点恰由既在极化根基内、又满足$q(v)=0$的非零向量给出（过代数闭包检验）。若极化非退化，这样的向量当然不存在。

例如特征$2$下$x^2+yz$的极化根基是$x$轴，$q$在其非零向量上不为零。根基的射影点$[1:0:0]$不在射影二次超曲面上，后者光滑；仿射锥原点却是奇点，因为三个偏导在那里均为零。这是作者所说“普通二次”与“极化非退化”不可混同的例子。

\paragraph{何时可逆推。}
在剩余特征非$2$时，$b_q(v,v)=2q(v)$，故极化一旦有非零根基，就产生几何奇点。在特征$2$且$E$秩为偶数时，极化为交替形式，秩为偶数；一旦退化，根基维数至少为二。在代数闭包上，$q$限制到根基后是一个线性形式的平方，故在根基中存在非零的零向量，同样产生几何奇点。因而在这两类情形，射影零点光滑可反推极化非退化。

\paragraph{术语的强弱。}
超奇异的定义是“不存在一个生成整个初始关系理想的非退化二次型”；可能根本没有单个二次生成元，并非仅有一个退化二次型生成理想。极超奇异的子预概形及修正判别式是本段提出的安排；本批不把这项提议写成原稿已经完成的独立定理。
\endgroup
